% Options for packages loaded elsewhere
\PassOptionsToPackage{unicode}{hyperref}
\PassOptionsToPackage{hyphens}{url}
\documentclass[
]{article}
\usepackage{xcolor}
\usepackage{amsmath,amssymb}
\setcounter{secnumdepth}{-\maxdimen} % remove section numbering
\usepackage{iftex}
\ifPDFTeX
  \usepackage[T1]{fontenc}
  \usepackage[utf8]{inputenc}
  \usepackage{textcomp} % provide euro and other symbols
\else % if luatex or xetex
  \usepackage{unicode-math} % this also loads fontspec
  \defaultfontfeatures{Scale=MatchLowercase}
  \defaultfontfeatures[\rmfamily]{Ligatures=TeX,Scale=1}
\fi
\usepackage{lmodern}
\ifPDFTeX\else
  % xetex/luatex font selection
\fi
% Use upquote if available, for straight quotes in verbatim environments
\IfFileExists{upquote.sty}{\usepackage{upquote}}{}
\IfFileExists{microtype.sty}{% use microtype if available
  \usepackage[]{microtype}
  \UseMicrotypeSet[protrusion]{basicmath} % disable protrusion for tt fonts
}{}
\makeatletter
\@ifundefined{KOMAClassName}{% if non-KOMA class
  \IfFileExists{parskip.sty}{%
    \usepackage{parskip}
  }{% else
    \setlength{\parindent}{0pt}
    \setlength{\parskip}{6pt plus 2pt minus 1pt}}
}{% if KOMA class
  \KOMAoptions{parskip=half}}
\makeatother
\usepackage{color}
\usepackage{fancyvrb}
\newcommand{\VerbBar}{|}
\newcommand{\VERB}{\Verb[commandchars=\\\{\}]}
\DefineVerbatimEnvironment{Highlighting}{Verbatim}{commandchars=\\\{\}}
% Add ',fontsize=\small' for more characters per line
\newenvironment{Shaded}{}{}
\newcommand{\AlertTok}[1]{\textcolor[rgb]{1.00,0.00,0.00}{\textbf{#1}}}
\newcommand{\AnnotationTok}[1]{\textcolor[rgb]{0.38,0.63,0.69}{\textbf{\textit{#1}}}}
\newcommand{\AttributeTok}[1]{\textcolor[rgb]{0.49,0.56,0.16}{#1}}
\newcommand{\BaseNTok}[1]{\textcolor[rgb]{0.25,0.63,0.44}{#1}}
\newcommand{\BuiltInTok}[1]{\textcolor[rgb]{0.00,0.50,0.00}{#1}}
\newcommand{\CharTok}[1]{\textcolor[rgb]{0.25,0.44,0.63}{#1}}
\newcommand{\CommentTok}[1]{\textcolor[rgb]{0.38,0.63,0.69}{\textit{#1}}}
\newcommand{\CommentVarTok}[1]{\textcolor[rgb]{0.38,0.63,0.69}{\textbf{\textit{#1}}}}
\newcommand{\ConstantTok}[1]{\textcolor[rgb]{0.53,0.00,0.00}{#1}}
\newcommand{\ControlFlowTok}[1]{\textcolor[rgb]{0.00,0.44,0.13}{\textbf{#1}}}
\newcommand{\DataTypeTok}[1]{\textcolor[rgb]{0.56,0.13,0.00}{#1}}
\newcommand{\DecValTok}[1]{\textcolor[rgb]{0.25,0.63,0.44}{#1}}
\newcommand{\DocumentationTok}[1]{\textcolor[rgb]{0.73,0.13,0.13}{\textit{#1}}}
\newcommand{\ErrorTok}[1]{\textcolor[rgb]{1.00,0.00,0.00}{\textbf{#1}}}
\newcommand{\ExtensionTok}[1]{#1}
\newcommand{\FloatTok}[1]{\textcolor[rgb]{0.25,0.63,0.44}{#1}}
\newcommand{\FunctionTok}[1]{\textcolor[rgb]{0.02,0.16,0.49}{#1}}
\newcommand{\ImportTok}[1]{\textcolor[rgb]{0.00,0.50,0.00}{\textbf{#1}}}
\newcommand{\InformationTok}[1]{\textcolor[rgb]{0.38,0.63,0.69}{\textbf{\textit{#1}}}}
\newcommand{\KeywordTok}[1]{\textcolor[rgb]{0.00,0.44,0.13}{\textbf{#1}}}
\newcommand{\NormalTok}[1]{#1}
\newcommand{\OperatorTok}[1]{\textcolor[rgb]{0.40,0.40,0.40}{#1}}
\newcommand{\OtherTok}[1]{\textcolor[rgb]{0.00,0.44,0.13}{#1}}
\newcommand{\PreprocessorTok}[1]{\textcolor[rgb]{0.74,0.48,0.00}{#1}}
\newcommand{\RegionMarkerTok}[1]{#1}
\newcommand{\SpecialCharTok}[1]{\textcolor[rgb]{0.25,0.44,0.63}{#1}}
\newcommand{\SpecialStringTok}[1]{\textcolor[rgb]{0.73,0.40,0.53}{#1}}
\newcommand{\StringTok}[1]{\textcolor[rgb]{0.25,0.44,0.63}{#1}}
\newcommand{\VariableTok}[1]{\textcolor[rgb]{0.10,0.09,0.49}{#1}}
\newcommand{\VerbatimStringTok}[1]{\textcolor[rgb]{0.25,0.44,0.63}{#1}}
\newcommand{\WarningTok}[1]{\textcolor[rgb]{0.38,0.63,0.69}{\textbf{\textit{#1}}}}
\usepackage{longtable,booktabs,array}
\usepackage{caption}
\captionsetup[table]{skip=6pt}
\newcounter{none} % for unnumbered tables
\usepackage{calc} % for calculating minipage widths
% Correct order of tables after \paragraph or \subparagraph
\usepackage{etoolbox}
\makeatletter
\patchcmd\longtable{\par}{\if@noskipsec\mbox{}\fi\par}{}{}
\makeatother
\setlength{\emergencystretch}{3em} % prevent overfull lines
\providecommand{\tightlist}{%
  \setlength{\itemsep}{0pt}\setlength{\parskip}{0pt}}
\usepackage{bookmark}
\IfFileExists{xurl.sty}{\usepackage{xurl}}{} % add URL line breaks if available
\urlstyle{same}
% fallback for those not using the hyperref driver hyperxmp:
\makeatletter
\@ifundefined{xmpquote}{\newcommand{\xmpquote}[1]{#1}}{}
\makeatother
\hypersetup{
  hidelinks,
  pdfcreator={LaTeX via pandoc}}

\author{}
\date{}

\begin{document}

\section{Cubli Control}\label{cubli-control}

Build a reaction-wheel balancing cube and compare PD/LQR with
Incremental Nonlinear Dynamic Inversion (INDI).

\subsection{Start here}\label{start-here}

\begin{enumerate}
\def\labelenumi{\arabic{enumi}.}
\tightlist
\item
  Solve and simulate a \textbf{one-wheel inverted pendulum} before
  designing the cube.
\item
  Size the motor and wheel together; CAD the frame, mounts, and wheel
  guard.
\item
  Build one restrained axis; verify IMU, motor torque control, and speed
  limits.
\item
  Establish PD/LQR balancing, then test INDI under the same
  disturbances.
\item
  Expand to three orthogonal wheels. Leave jump-up for a later stage.
\end{enumerate}

Print brackets and fit-check parts, not an unqualified high-speed
flywheel. Use a secured metal rotor, containment, and an emergency stop.

\subsection{Build the documents}\label{build-the-documents}

Pandoc and MiKTeX are used to turn the Markdown documents into LaTeX and
PDFs. Run the VS Code task \textbf{Build Markdown PDFs} to create
\texttt{build/latex/README.tex} and \texttt{build/latex/STUDY.tex}, plus
\texttt{README.pdf} and \texttt{Ref/STUDY.pdf}. The \textbf{Watch
Markdown PDFs} task builds both documents and rebuilds the affected PDF
when either Markdown file is saved. Missing MiKTeX packages are
installed automatically during compilation.

\subsection{1D simulation}\label{d-simulation}

Start with one reaction wheel and a fixed, frictionless edge pivot. The
parameters in \href{Matlab\%20sim/Params.m}{Params.m} are the CAD values
from Table I of
\href{Ref/lagrange\%20cubli\%20\%26\%20complex\%20formation.pdf}{Bobrow
et al.}, not measured prototype values.

\subsubsection{Coordinates and the paper's missing
offset}\label{coordinates-and-the-papers-missing-offset}

The simulation now uses the paper's Fig. 4 coordinates:

\begin{itemize}
\tightlist
\item
  \(\theta_c=\) \texttt{thetab}: angle from world horizontal \(+Y\) to a
  cube face.
\item
  \(\theta_w=\) \texttt{thetaw}: reaction-wheel angle relative to the
  cube.
\item
  \(\omega_w=\dot\theta_w=\) \texttt{thetawdot}: relative wheel speed.
\item
  The center-of-mass line is at \(\theta_c+\pi/4\).
\item
  The unstable edge-balancing position is \(\theta_c=\pi/4=45^\circ\),
  not zero.
\item
  Positive motor torque accelerates the wheel and applies negative
  reaction torque to the body.
\end{itemize}

The paper gives

\[T=\frac12\bar I_{cO}\dot\theta_c^2+
  \frac12 I_{wG}(\dot\theta_c+\dot\theta_w)^2,
\qquad
V=m_cgd\sin\left(\theta_c+\frac{\pi}{4}\right).\]

Therefore,

\[\frac{\partial V}{\partial\theta_c}
=m_cgd\cos\left(\theta_c+\frac{\pi}{4}\right).\]

The printed Eq. (7), Eq. (9), Eq. (11), and Fig. 5 show
\(m_cgd\cos\theta_c\) without the \(\pi/4\). That expression conflicts
with the paper's Eq. (6), Fig. 4 geometry, and its later statement that
\(\theta_c=45^\circ\) is the unstable equilibrium. The simulation
retains the paper coordinate and cosine gravity term, but restores the
required offset:

\[\bar I_{cO}\ddot\theta_c+
I_{wG}(\ddot\theta_c+\ddot\theta_w)
=-m_cgd\cos\left(\theta_c+\frac{\pi}{4}\right),\]

\[I_{wG}(\ddot\theta_c+\ddot\theta_w)
=-\tau_f(\dot\theta_w)+\tau.\]

Let \(u=\tau-\tau_f(\dot\theta_w)\). The exact matrix model
corresponding to the paper's Eq. (9) is

\[\begin{bmatrix}
\bar I_{cO} & I_{wG}\\
0 & I_{wG}
\end{bmatrix}
\begin{bmatrix}
\ddot\theta_c\\
\ddot\theta_c+\ddot\theta_w
\end{bmatrix}
=
\begin{bmatrix}
-m_cgd\cos(\theta_c+\pi/4)\\
u
\end{bmatrix}.\]

Solving it gives

\[\ddot\theta_c=
\frac{-m_cgd\cos(\theta_c+\pi/4)-u}{\bar I_{cO}},
\qquad
\ddot\theta_w=\frac{u}{I_{wG}}-\ddot\theta_c.\]

The second equation matters because \(\theta_w\) is relative to the
body. Paper Eq. (11) uses \(\ddot\theta_w\gg\ddot\theta_c\) to
approximate \(\ddot\theta_w\approx u/I_{wG}\); the MATLAB script instead
integrates the exact Eq. (9) dynamics. Motor friction uses paper Eq.
(36) with the identified Table II values:

\[\tau_f(\dot\theta_w)=\operatorname{sign}(\dot\theta_w)
\left(\tau_c+b_w|\dot\theta_w|+c_d\dot\theta_w^2\right),\]

\(\tau_c=2.46\times10^{-3}\) N m, \(b_w=1.06\times10^{-5}\) N m s/rad,
\(c_d=1.70\times10^{-8}\) N m s²/rad². \texttt{sign} is smoothed as
\(\tanh(\dot\theta_w/\omega_\epsilon)\), \(\omega_\epsilon=0.5\) rad/s,
to avoid solver chattering at zero wheel speed. Optional pivot damping
\(-b_c\dot\theta_c\) is added to the body equation; the paper has none,
so \texttt{bc\ =\ 0} by default.

\subsubsection{Why the swinging cube barely loses
energy}\label{why-the-swinging-cube-barely-loses-energy}

Motor friction is \textbf{internal}: it acts between wheel and body, and
it only dissipates power \(\tau_f\dot\theta_w\) when the wheel slips
relative to the body. The paper has no pivot/edge friction, so the
swinging body itself is conservative. With a small swing about the
hanging position, the validated model loses about 15\% of the swing
energy in 10 s. Set \texttt{bc\ \textgreater{}\ 0} to model edge
friction if your hardware decays faster.

In paper Eq. (11)/Fig. 5, the relative wheel acceleration is
approximated as \((\tau-\tau_f)/I_{wG}\), independent of body motion. In
that approximation, with \(\tau=0\) and \(\dot\theta_w(0)=0\), the wheel
never slips, \(\tau_f=0\), and the body loses \textbf{no} energy at all.

\subsubsection{\texorpdfstring{Why constant torque barely moves the body
with large
\texttt{bw}}{Why constant torque barely moves the body with large bw}}\label{why-constant-torque-barely-moves-the-body-with-large-bw}

The body only feels the motor reaction while the wheel accelerates:
\(I_{wG}(\ddot\theta_c+\ddot\theta_w)=\tau-\tau_f\). Once friction
equals the motor torque, \(u=0\) and the body sees no net torque. With a
viscous \texttt{bw\ =\ 1e-2}, the wheel reaches \(\tau/b_w\) (20 rad/s
for 0.2 N m) in about \(I_{wG}/b_w=12.5\) ms, giving the body only a
tiny impulse: \(\theta_b\) moved about \(\pm1.3^\circ\). With the
paper's friction, the same torque spins the wheel to about 3000 rad/s
over several seconds, and \(\theta_b\) swings from \(-115^\circ\) to
\(-77^\circ\). Also note gravity torque is \(mgd\approx0.88\) N m, so a
constant hanging-position offset needs \(\sin^{-1}(\tau/mgd)\); that is
only \(3.2^\circ\) for 0.05 N m.

\subsubsection{Validation results}\label{validation-results}

{\def\LTcaptype{none} % do not increment counter
\begin{longtable}[]{@{}
  >{\raggedright\arraybackslash}p{(\linewidth - 2\tabcolsep) * \real{0.5000}}
  >{\raggedright\arraybackslash}p{(\linewidth - 2\tabcolsep) * \real{0.5000}}@{}}
\toprule\noalign{}
\begin{minipage}[b]{\linewidth}\raggedright
Check
\end{minipage} & \begin{minipage}[b]{\linewidth}\raggedright
Result
\end{minipage} \\
\midrule\noalign{}
\endhead
\bottomrule\noalign{}
\endlastfoot
Simulink vs \texttt{ode45}, 80° release, 0--2 s &
\(\max|\Delta\theta_b|=6\times10^{-4}\) rad \\
Small-swing period about hanging & 0.775 s vs theory
\(2\pi/\sqrt{mgd/\bar I_{cO}}=0.771\) s \\
Hanging, 0.05 N m & body settles about 2.9° from hanging, near
\(\sin^{-1}(\tau/mgd)=3.2^\circ\) \\
Upright 90°, zero torque & stays briefly, then falls (unstable
equilibrium) \\
\end{longtable}
}

\subsubsection{Set the initial
conditions}\label{set-the-initial-conditions}

Edit these values in \href{Matlab\%20sim/Params.m}{Params.m}:

\begin{Shaded}
\begin{Highlighting}[]
\VariableTok{thetab0} \OperatorTok{=} \VariableTok{deg2rad}\NormalTok{(}\FloatTok{90}\NormalTok{)}\OperatorTok{;} \CommentTok{\% theta\_c; upright is 90 deg when com\_angle\_offset = 0}
\VariableTok{thetabdot0} \OperatorTok{=} \FloatTok{0}\OperatorTok{;}        \CommentTok{\% body speed, rad/s}
\VariableTok{thetaw0} \OperatorTok{=} \FloatTok{0}\OperatorTok{;}           \CommentTok{\% wheel angle relative to body, rad}
\VariableTok{thetawdot0} \OperatorTok{=} \FloatTok{0}\OperatorTok{;}        \CommentTok{\% wheel speed relative to body, rad/s}
\end{Highlighting}
\end{Shaded}

The Simulink integrators read these names directly; never type numbers
such as \texttt{pi} into an integrator's initial condition. Then
\texttt{x0\ =\ {[}thetab0;\ thetabdot0;\ thetaw0;\ thetawdot0{]}}.
Upright is \texttt{pi/2\ -\ com\_angle\_offset} and hanging is
\texttt{-pi/2\ -\ com\_angle\_offset}. To observe falling, start a few
degrees away from upright.

\subsubsection{Simulink blocks}\label{simulink-blocks}

Run \texttt{Params} in the model's InitFcn callback (script must be on
the MATLAB path). All angles entering trigonometric blocks are radians.
Do not multiply an integrator output by \texttt{pi/180}; use
\texttt{deg2rad} only for numeric initial values.

\begin{enumerate}
\def\labelenumi{\arabic{enumi}.}
\tightlist
\item
  Subtract body speed from absolute wheel speed:
  \(\dot\phi=s-\dot\theta\). In the four-state implementation this
  signal is directly \texttt{thetawdot}.
\item
  The \texttt{Friction} Fcn block computes paper Eq. (36) from
  \texttt{thetawdot}; subtract it from motor torque to obtain \(u\).
\item
  Add \texttt{com\_angle\_offset} to \texttt{thetab} (\texttt{ComOffset}
  Bias block), apply \texttt{cos}, then multiply by
  \texttt{gravity\_gain}. The body sum \texttt{Add} is
  \(\tau_f-\tau-mgd\cos(\cdot)-b_c\dot\theta_c\); apply
  \texttt{body\_accel\_gain} to obtain \(\ddot\theta_c\).
\item
  Integrate \(\ddot\theta_c\) to \(\dot\theta_c\), then to \(\theta_c\).
  Initial conditions: \texttt{thetabdot0}, then \texttt{thetab0}.
\item
  Apply \texttt{wheel\_accel\_gain} to \(u\) to obtain absolute wheel
  acceleration \(\ddot\theta_c+\ddot\theta_w\). Subtract
  \(\ddot\theta_c\) to obtain relative \(\ddot\theta_w\), then integrate
  with \texttt{thetawdot0} and \texttt{thetaw0}.
\end{enumerate}

The original pictured model used \texttt{cos(thetab*pi/180)}, making
gravity nearly maximum at its zero initial state. With the listed
values, \texttt{mgd} is about 0.8844 N m and \(\bar I_{cO}\) is
0.0133125 kg m², so a 0.19 N m input at zero produces a very large
acceleration. That explains the sharp scope drop; it is not an
integration discontinuity.

\subsubsection{MATLAB simulation and 3D
animation}\label{matlab-simulation-and-3d-animation}

Run \href{Matlab\%20sim/cubli_3d_compare.m}{cubli\_3d\_compare.m}. It
performs these steps:

\begin{enumerate}
\def\labelenumi{\arabic{enumi}.}
\tightlist
\item
  Runs \texttt{Params.m}.
\item
  Runs the Simulink model with \texttt{RelTol\ =\ 1e-8} and
  \texttt{MaxStep\ =\ dt}, then reads the \textbf{To Workspace} outputs
  \texttt{thetab}, \texttt{thetabddot}, \texttt{thetawdot}, and
  optionally \texttt{thetaw} (otherwise integrated from
  \texttt{thetawdot}).
\item
  Runs the exact nonlinear Eq. (9) model with \texttt{ode45} as a
  reference.
\item
  Resamples both on a uniform time grid and stores the Simulink result
  in the workspace table \texttt{measurements}.
\item
  Overlays Simulink and \texttt{ode45} for \(\theta_b\),
  \(\ddot\theta_b\), relative \(\dot\theta_w\), and total energy
  \(\tfrac12\bar I_{cO}\dot\theta_c^2+\tfrac12 I_{wG}(\dot\theta_c+\dot\theta_w)^2+mgd\sin(\theta_c+\delta)\).
\item
  Animates the selected source (\texttt{animationSource\ =\ "simulink"}
  or \texttt{"ode"}) about the fixed \(X\) edge. Red marks the pivot,
  RGB lines are body axes, magenta is the wheel axis, and the orange
  spoke shows wheel rotation.
\end{enumerate}

\subsubsection{Sending Simulink output to the
animation}\label{sending-simulink-output-to-the-animation}

The model
\href{Matlab\%20sim/dynamic_of_reaction_wheel.slx}{dynamic\_of\_reaction\_wheel.slx}
already contains these \textbf{To Workspace} blocks (Save format
\texttt{Timeseries}), each connected to a \textbf{radian} signal, not to
a \texttt{180/pi} scope gain:

{\def\LTcaptype{none} % do not increment counter
\begin{longtable}[]{@{}
  >{\raggedright\arraybackslash}p{(\linewidth - 2\tabcolsep) * \real{0.5000}}
  >{\raggedright\arraybackslash}p{(\linewidth - 2\tabcolsep) * \real{0.5000}}@{}}
\toprule\noalign{}
\begin{minipage}[b]{\linewidth}\raggedright
Variable name
\end{minipage} & \begin{minipage}[b]{\linewidth}\raggedright
Connect to
\end{minipage} \\
\midrule\noalign{}
\endhead
\bottomrule\noalign{}
\endlastfoot
\texttt{thetab} & \texttt{Integrator1} (second body integrator),
\(\theta_b\) \\
\texttt{thetabdot} & \texttt{Integrator} (first body integrator),
\(\dot\theta_b\) \\
\texttt{thetabddot} & \texttt{Gain} (\texttt{body\_accel\_gain}),
\(\ddot\theta_b\) \\
\texttt{thetawdot} & \texttt{Integrator3} (first wheel integrator),
relative \(\dot\theta_w\) \\
\texttt{thetaw} & \texttt{Integrator2} (second wheel integrator) \\
\end{longtable}
}

For \texttt{thetawdot} to be relative to the body, the wheel branch is:

\begin{enumerate}
\def\labelenumi{\arabic{enumi}.}
\tightlist
\item
  \texttt{Add1\ =\ tau\ -\ tau\_f} (second input from the
  \texttt{Friction} block, \textbf{not} from the body sum).
\item
  \texttt{wheel\_accel\_gain} gives absolute acceleration
  \(\ddot\theta_b+\ddot\theta_w=u/I_{wG}\).
\item
  \texttt{Add2} (\texttt{+-}) subtracts the \texttt{body\_accel\_gain}
  \textbf{output} \(\ddot\theta_b\) (an acceleration, not the torque sum
  \texttt{Add}), giving relative \(\ddot\theta_w\) for the wheel-speed
  integrator.
\end{enumerate}

Integrator initial conditions are \texttt{thetabdot0}, \texttt{thetab0},
\texttt{thetawdot0}, and \texttt{thetaw0}. The gravity block uses
\texttt{cos(thetab\ +\ com\_angle\_offset)}; the animation draws the
center of mass at the same angle:

\[\mathbf r_G=
\begin{bmatrix}
0\\d\cos(\theta_b+\delta)\\d\sin(\theta_b+\delta)
\end{bmatrix},
\qquad \delta=\texttt{com\_angle\_offset}.\]

With \(\delta=0\) (printed Eq. (7)), upright is \(\theta_b=90^\circ\);
with \(\delta=\pi/4\) (Eq. (6)), upright is \(45^\circ\).

There is no floor collision, hard stop, motor saturation, or feedback
controller in this open-loop model. Once the cube falls past the
physical floor, the mathematical model continues rotating; add contact
dynamics separately if impacts are required.

\subsection{Study material}\label{study-material}

\begin{itemize}
\tightlist
\item
  \href{Ref/STUDY.md}{ETH thesis summary, textbook reading, and math
  exercise}
\item
  \href{Ref/phdthesisE-Collection.pdf}{Attached ETH thesis: Michael
  Muehlebach, 2018}
\item
  \href{https://idsc.ethz.ch/research-dandrea/research-projects/archive/cubli.html}{ETH
  Cubli project and papers}
\item
  \href{https://www.egr.msu.edu/~khalil/NonlinearSystems/index.html}{Khalil,
  \emph{Nonlinear Systems}, 3rd edition: author's companion site}
\end{itemize}

\end{document}
